\section{Preliminaries}
\label{sec:preliminary}

% Iterate multiple times when generating a single token
\xhdr{Autoregressive LLMs}
Modern LLMs generate text through an autoregressive next-token prediction process.
It includes a \emph{prefill} stage and a \emph{decode} stage~\citep{radford2018gpt1,radford2019gpt2,kwon2023vllm}.
In the prefill stage, the model processes the entire input sequence in parallel; in the decode stage, it consumes one new token at a time along with cached history to predict the next token.

Formally, let $t_i$ denote the token at position $i$ and $x_i\in\mathbb{R}^h$ its embedding. 
Let $E\in\mathbb{R}^{v\times h}$ be the embedding matrix, so $x_i=E[t_i]$ when $t_i$ is treated as an index. 
Here, $v$ and $h$ are the vocabulary size and hidden dimension. 
The output projection matrix is $W_{\text{out}}\in\mathbb{R}^{h\times v}$ (equal to $E^\top$ if tied). 
Given the context $T_{\le i}=[t_0,\dots,t_i]$ with embeddings $X_{\le i}=[x_0,\dots,x_i]$, the model $\theta$ produces a \emph{last-layer hidden state} $y_i$ for token $t_i$:
\begin{equation}
    y_i \;=\; \mathcal{P}_{\theta}\!\big(x_i \mid X_{\le i}\big)\;\in\;\mathbb{R}^h .
    \label{eq:llm_decode}
\end{equation}
The next-token distribution $p_i$ and sample are:
\begin{equation}
    p_i \;=\; \softmax\!\big(W_{\text{out}}^\top y_i\big) \in \mathbb{R}^v,
    \qquad
    t_{i+1} \;=\; \mathcal{S}(p_i),
    \label{eq:llm_sample}
\end{equation}
where $\mathcal{S}$ is a sampling rule such as nucleus sampling. 
Decoding repeats until an end-of-sequence token is generated.

\xhdr{Causal attention}
Modern LLMs typically adopt a \emph{causal} attention mechanism.
As shown in Figure~\ref{fig:arch}(a), each position attends only to itself and earlier positions, consistent with Equation~\ref{eq:llm_decode}. 
This design brings two key benefits: 
(1) it enables parallel training with next-token prediction and shifted logits, avoiding the need for token-by-token generation; 
and (2) it allows efficient inference by caching Key/Value states of past tokens instead of recomputing them.

\xhdr{Looped transformers}
Looped transformers introduce an inner loop that iterates in latent space before verbalizing each output token. 
Let $d\in\{1,2,\dots\}$ denote the iteration depth (written as a superscript), and set $x_i^{(0)}=E[t_i]$. 
At each iteration, looped transformers update $y_i$ with causal attention on the hidden states of \emph{the current iteration}:
\begin{equation}
    y_i^{(d)} \;=\; \mathcal{P}_{\theta}\!\big(x_i^{(d)} \,\big|\, X_{\le i}^{(d)}\big),
    \qquad
    X_{\le i}^{(d)}=[x_0^{(d)},\dots,x_i^{(d)}].
    \label{eq:rllm_decode}
\end{equation}
An inner transition then produces the next-depth embedding. 
For example, Loop~\citep{saunshi2025loopedTrans} simply sets $x_i^{(d+1)}=y_i^{(d)}$, while Ponder~\citep{zeng2025pondering} uses logit-weighted embeddings:
\begin{equation}
    x_i^{(d+1)}=\softmax\!\big(W_{\text{out}}^\top y_i^{(d)}\big)\ E = p_i^{(d)} E.
    \label{eq:ours_update}
\end{equation}
In practice, it uses the top-100 logits instead of full logits for efficiency.
% looped transformer D
% ponder D; logits
% MoR D or 0

Verbalization occurs at a fixed \emph{maximum depth} $d_{\max}$ shared by all tokens, where $y_i^{(d_{\max})}$ is transformed into the next token $t_{i+1}$, resembling Equation~\ref{eq:llm_sample}.

%%%%%%%%%%%%%%%%%%% Latent Overthinking %%%%%%%%%%%%%%%%%%%%
\section{Selective Latent Iteration Oracles}
\label{sec:oracle}

\begin{table}[t]
    \centering
    \footnotesize
    \setlength\tabcolsep{2pt}
    \caption{Performance comparison across iteration strategies, using Ouro-1.7B model except for the last row. Always1 means no latent iteration. Oracle policy iterates on 12-19\% of tokens. MMLU100 denotes the first 100 questions from MMLU-STEM.}
    \definecolor{tahgray}{RGB}{242,242,242}
    \begin{tabular}{l@{\hskip 3pt}cccc}
    \toprule
    \textbf{Iter. Policy} & \textbf{NTP} & \textbf{AMC23} & \textbf{MMLU100} & \textbf{HE++} \\
    \midrule
    % \multicolumn{5}{c}{Ouro Model}\\
    % \midrule
    Always1 & 73.1 & 38.1 & 56.0 & 39.6 \\
    Always2 & 79.7 & 40.6 & 60.0 & 40.9 \\
    \rowcolor{tahgray}
    Orcl.  &
    81.8{\scriptsize$/+2.1$} &
    47.9{\scriptsize$/+7.3$} &
    62.0{\scriptsize$/+2.0$}      &
    43.3{\scriptsize$/+2.4$} \\
    \midrule
    % \multicolumn{5}{c}{\name Model}\\
    % \midrule
    % Always1 & 78.7 & 40.9 & 60.0 & 45.4 \\
    % Always2 & 47.6 & 1.3  & 20.0 & 0.0  \\
    \rowcolor{tahgray}
    Orcl. w.\name &
    89.3{\scriptsize$/+9.6$} & 
    68.8{\scriptsize$/+28.2$} & 
    85.0{\scriptsize$/+25.0$} & 
    72.9{\scriptsize$/+32.0$}
          \\
    \bottomrule
    \end{tabular}
    \label{tab:oracle_policy}
\end{table}

\begin{figure*}[t]
    \centering
    \includegraphics[width=\textwidth]{figure/arch.pdf}
    \caption{\name Overview. 
    (a) Regular causal attention: tokens attend only to previous positions.
    (b) Our duo-causal attention: tokens attend to both previous positions and shallower iteration depths, maintaining 2D causality.
    (c) Model architecture: \name selectively iterates or verbalizes tokens. It uses LoRA at deeper iterations to shift from next-token prediction to hard-token refinement. A neural decider determines whether to continue iterating or output the token.
    }
    \label{fig:arch}
\end{figure*}

\xhdr{Setup}
We analyze the latent iteration behavior of the Ouro looped transformer~\citep{zhu2025ouro} and our proposed \name model (architecture detailed in Section~\ref{sec:method}). 
All models are fine-tuned from Qwen3-1.7B-Base on a balanced 100K-sample subset of the Open-R1 dataset~\citep{openr1}, following the setups in Section~\ref{sec:eval_setup}.

\xhdr{Oracle iteration policy}
To investigate the potential of selective iteration, we establish an oracle policy $\pi$. It triggers additional iterations only when the reference LLM $\theta$ mispredicts the target token at the first forward pass.
In this section, we set $\theta$ to the model under evaluation, so $\pi$ acts as a greedy, locally optimal policy.

Formally, let $\hat{p}_i$ denote the reference model's first-pass next-token distribution at position $i$, and $t_{i+1}$ the ground-truth token.
The oracle iteration depth $d_i^{\pi}$ is:
\begin{equation}
    d_i^{\pi} \;=\; 1+\mathcal{D}(\hat{p}_i^{(1)}, t_{i+1}),
    \label{eq:ours_oracle_depth}
\end{equation}
where $\mathcal{D}$ is a binary discrepancy metric.
We use top-1 mismatch here: $\mathcal{D} = \1[\hat{t}_{i+1} \neq t_{i+1}]$, with $\hat{t}_{i+1}=\arg\max_{t} \hat{p}_i^{(1)}$ ($\1[\cdot]$ denotes the indicator function).
For simplicity, we assume $d_{\max}=2$. The arbitrary-depth case is in Appendix~\ref{sec:appendix/oracle_policy}, and alternative discrepancy metrics are ablated in Table~\ref{tab:ablation}.

For brevity, we call a token \emph{easy} if the reference model correctly predicts it at the first pass ($d_i^{\pi}=1$), and \emph{hard} otherwise.
Prior work~\citep{fu2025r2r} shows that \emph{hard} tokens typically occupy only a small proportion (e.g., $7\%$) of all tokens.

\xhdr{Next-token prediction}
We evaluate next-token prediction (NTP) accuracy by comparing top-1 predictions against ground-truth tokens in the Open-R1 validation set.
As shown in Table~\ref{tab:oracle_policy}, always iterating twice improves NTP accuracy for Ouro by 6.6\%, but at the cost of doubled depth.
Surprisingly, the oracle policy, which skips iterations on 81-88\% of tokens, further improves accuracy by 2.1\%.
This gain comes from avoiding \emph{latent overthinking}: without selective iteration, the model can revise correct predictions to wrong ones at the second pass, as shown in Figure~\ref{fig:landscape}.

\xhdr{Downstream tasks}
We next examine whether NTP gains translate to downstream task performance.
Since ground-truth tokens are unavailable during generation, we use top-1 predictions from the stronger Qwen3-8B as proxy labels.
Table~\ref{tab:oracle_policy} shows that the oracle policy improves downstream performance by 2.0-7.3\%.

\xhdr{\name design objectives}
The oracle experiments reveal two key insights.
First, models have significant untapped potential when iteration depth is selected correctly.
Second, the oracle policy is effective enough to learn from, even though other globally optimal policies may exist.
These findings motivate the design objectives of \name: (1) its model architecture should natively support selective iteration depths; and (2) following the oracle policy, its training objective for deeper iterations is not to cover all tokens, but to selectively focus only on the few failing tokens.
We later validate that \name-1.7B can better utilize the oracle policy to achieve $>25\%$ improvement, surpassing even Qwen3-4B.
While the oracle requires ground-truth tokens unavailable at inference, approximating it via neural networks is a promising direction.

%%%%%%%%%%%%%%%%%%% Our Design %%%%%%%%%%%%%%%%%%%%
\section{\name Design}
\label{sec:method}

We expand the key motivations and designs of \name in this section, including the duo-causal attention mechanism for cross-depth information flow (Section~\ref{sec:duo_attention}), the depth-adaptive model architecture with LoRA adapters (Section~\ref{sec:model_architecture}), and a two-stage training scheme (Section~\ref{sec:training}).

\subsection{Duo-Causal Attention}
\label{sec:duo_attention}
% Dynamic latent thinking poses a unique challenge: how to maximize information flow across tokens with diverse iteration depths while maintaining training and prefill parallelism. We address this with a duo-causal attention mechanism.
\xhdr{Motivation}
In looped transformers with fixed depth, standard causal attention on the current iteration's KV states incorporates all context (Equation~\ref{eq:rllm_decode}). 
However, dynamic iteration depths create a challenge: tokens at deeper levels cannot access hidden states of previous tokens that verbalized at shallower depths.
This poses a dilemma between requiring up-to-date context from all previous tokens and maintaining parallel training where depth-$d$ computations cannot depend on uncomputed deeper states ($d'>d$).
Existing approaches compromise on one of these aspects. Some sacrifice parallelism by allowing attention to deeper iterations~\citep{hao2024training}; others preserve parallelism by restricting attention to only the initial iteration's KVs~\citep{bae2025mor}.
To resolve this dilemma, we introduce a simple yet effective mechanism to maximize cross-depth information flow while maintaining high parallelism.

\xhdr{Duo-causal attention mechanism}
As shown in Figure~\ref{fig:arch}(b), duo-causal attention extends \emph{causality} to two dimensions, letting tokens attend across both previous positions and shallower iteration depths.
Formally, we extend the accessible set from Equation~\ref{eq:rllm_decode} to
\begin{equation}
    X_{\le i}^{(\le d)} \;=\; \{\, x_j^{(k)} \;|\; j\le i,\; k\le d \,\}.
    \label{eq:duo_causal_attention}
\end{equation}
When all tokens iterate only once (as in standard transformers), this reduces to regular causal attention. 
The duo-causal design achieves both full parallel training and cross-depth information flow.
At depth $d$, all tokens compute their depth-$d$ representations simultaneously using \emph{only and all} information from depths $1$ through $d$.

\begin{figure}[tb]
    \centering
    \includegraphics[width=\linewidth]{figure/duo_attention_expand.pdf}
    \caption{Duo-causal attention implementation. (a) Conceptual \name example with dynamic iteration depths. Each cell denotes a token–depth pair \((\text{token id}, \text{iter depth})\). (b) Each iteration maintains its own KV cache. (c) KV caches from all iterations are concatenated into a 1D sequence and processed with standard attention under a duo-causal mask. The duo-causal mask is conceptually partitioned into blocks by iteration depth. The diagonal blocks use a standard causal mask, while off-diagonal blocks use reduced causal masks that enforce the duo-causal rules.}
    \label{fig:duo_attention_expand}
\end{figure}

Duo-causal attention is fully compatible with attention kernels like FlashAttention~\citep{dao2022flashattention,dao2023flashattention2,shah2024flashattention3} and sparse implementations~\citep{fu2025moa, zhang2025spargeattention}.
As shown in Figure~\ref{fig:duo_attention_expand}, we simply maintain separate KV caches per iteration depth and flatten the 2D (token, depth) grid into 1D by concatenating deeper KV caches after shallower ones.
Positional encodings are applied based solely on the original token index, invariant to iteration depth.
The duo-causal constraint is then enforced via a modified additive attention mask, requiring no custom CUDA kernels.
More details on duo-causal attention implementation are discussed in Appendix~\ref{sec:appendix/duo-causal-attention}.

% With duo-causal attention, each token maximizes its contextual information while preserving the efficiency benefits of parallel training. This is particularly crucial for \name, where different tokens may terminate at different depths, requiring flexible information flow across the iteration landscape.
% \todo{we can remove this for conciseness}


%%%%%%%%%%%%%% model architecture %%%%%%%%%%%%%%
\subsection{Depth-adaptive Model Architecture}
\label{sec:model_architecture}
% Key innovations:
% 1. add LoRA for d larger than 1.
% 3. adaptively exit with iteration decider
% The architecture of \name is designed to dynamically shift prediction objectives across iterations.
% Figure~\ref{fig:arch}(c) shows the overall architecture of \name, containing an iteration-aware backbone and a neural iteration decider.

\xhdr{Motivation}
Previous looped transformers typically use identical weights across all iterations.
However, we find that over 73\% of next-tokens are correctly predicted at the first iteration (Table~\ref{tab:oracle_policy}).
This suggests deeper iterations serve a different objective: they refine the first iteration's prediction rather than predicting further ahead to the next-next token.
This mirrors deep LLMs, where shallow layers predict next tokens for deeper layers to refine~\citep{belrose2023tunedLens, schuster2022confidentEarlyExit, bae2023fastEarlyExit}.
While deep LLMs naturally handle this shift through distinct parameters per depth, looped transformers must accommodate both objectives with shared weights, potentially limiting performance.
Moreover, fixed iteration depths can cause \emph{latent overthinking}, motivating our dynamic approach.

\xhdr{Backbone model}
To address the objective shift, we apply a LoRA adapter~\citep{hu2022lora} to the shared LLM backbone only for iterations $d>1$.
As shown in Figure~\ref{fig:arch}(c), this allows the base LLM to focus on latent embeddings, while the adapter handles the objective shift.
It preserves strong next-token prediction at $d=1$, alleviating interference from deeper iterations.
We also add residual connections across iterations to simplify the refinement and improve gradient flow.
Formally, at depth $d$, we compute
\begin{equation}
    y_i^{(d)} \;=\; \mathcal{P}_{\theta_d}\!\Big(x_i^{(d)} \,\Big|\, X_{\le i}^{(\le d)}\Big),
    \label{eq:ours_decode}
\end{equation}
with depth-specific parameters
\[
\theta_d=\theta \ \text{for } d=1,
\qquad
\theta_d=\theta+\Delta \ \text{for } d > 1,
\]
where $\theta$ and $\Delta$ denote the LLM and LoRA weights, respectively.
The next-iteration inputs use logit-weighted embeddings (Equation~\ref{eq:ours_update}); verbalization follows standard sampling (Equation~\ref{eq:llm_sample}).
Each $y_i^{(d)}$ either continues iterating or verbalizes according to the decider $\mathcal{I}_\phi$.

\xhdr{Iteration decider}
We use a lightweight MLP as the iteration decider $\mathcal{I}_\phi$ to determine whether each token should continue iterating or verbalize.
After each iteration, it concatenates the backbone's shallow, middle, and final hidden states to predict a continuation probability:
\[
\hat{c}_i^{(d)} \;=\; \mathcal{I}_\phi\!\big(h_i^{(d)}\big) \in [0,1].
\]
During inference, token $i$ verbalizes when $\hat{c}_i^{(d)}$ falls below threshold $c_{\text{threshold}}$ or reaches maximum depth $d_{\max}$.

%%%%%%%%%%%%%% training scheme %%%%%%%%%%%%%%
\subsection{Training Scheme}
\label{sec:training}

We adopt a two-stage training scheme. We first fine-tune the backbone model to support selective latent iteration, then train the iteration decider. Both stages are aligned to the same oracle iteration policy.

\xhdr{Motivation}
As shown in Figure~\ref{fig:arch}(c), the backbone LLM and the iteration decider are tightly coupled, making joint training unstable.
Specifically, the backbone’s prediction quality across iterations determines the optimal depth, while the decider controls the backbone’s KV cache and output depth. 
To stabilize training, we train the two components sequentially under a fixed oracle policy $\pi$ (validated in Section~\ref{sec:oracle}).

\paragraph{Oracle Iteration Policy $\pi$.}
Our goal is to iterate only on \emph{hard} tokens that a standard supervised fine-tuned (SFT) model would mispredict on the first pass.
Thus, we define $\pi$ with the SFT model as the reference: we trigger an additional iteration when its top-1 next-token prediction differs from the ground-truth token.
In principle, one could instead define $\pi$ using the current looped model itself (i.e., an on-policy oracle), but we find this to be empirically unstable (Section~\ref{sec:exp/dse}).

% \paragraph{Oracle Iteration Policy $\pi$.}
% Since the desired behavior of our looped LLM is to selectively iterate at \emph{hard} tokens that standard SFT-variant of the LLM mispredict. Thus, we use the SFT model as the reference model, iterating when its top-1 prediction mismatches training data. 
% Note that while directly referring the training looped LLM itself is also an option, it is empirically unstable, as detailed in Section~\ref{sec:exp/dse}.

\paragraph{Stage 1: Backbone supervision under $\pi$.}
We optimize the backbone LLM ($\theta$ and LoRA adapter $\Delta$) with $\pi$-guided iteration execution.
The loss is standard next-token prediction at the oracle-determined depth:
\[
\mathcal{L}_{\text{SFT}}(\theta,\Delta)
\;=\;
\sum_i \; -\log p_i^{(d_i^{\pi})}\!\big(t_{i+1}\big),
\]
where $p_i^{(d_i^{\pi})}$ is the next-token distribution at position $i$, depth $d_i^{\pi}$.
This preserves first-iteration accuracy for easy tokens while training deeper iterations to refine hard tokens.

\paragraph{Stage 2: Decider imitation under frozen backbone.}
We freeze the backbone model $(\theta,\Delta)$ and train the iteration decider $\phi$ to imitate the oracle policy's continuation decisions.
We minimize weighted binary cross-entropy:
\[
\mathcal{L}_{\text{dec}}(\phi)
\;=\;
-\sum_{i,d} w_d 
c_i^{(d)} \log \hat{c}_i^{(d)} + (1-c_i^{(d)})\log(1-\hat{c}_i^{(d)}),
\]
where the sum ranges over tokens $i$ and depths $d=1,\ldots,\min\{d_{\max}-1, d_i^{\pi}\}$. Here $c_i^{(d)}$ is the ground-truth continuation label, $\hat{c}_i^{(d)}$ is the predicted probability, and $w_d$ is the class weight for label imbalance (ratio of stop to continue labels).

% This reweighting ensures stable training with unbalanced continue and stopping label.

% \paragraph{Inference.}
% At test time, for each token $i$ we iterate depths $d=0,1,\dots$ and stop when $q_i^{(d)} \le \tau$ or $d=D_{\max}$, then verbalize with $p_i^{(d)}$.
% The threshold $\tau$ can be tuned on a validation set to meet an accuracy–compute budget.
% This realizes dynamic latent thinking while preserving prefill parallelism via duo-causal attention.

The two-stage scheme stabilizes training by decoupling backbone learning (conditioned on a fixed $\pi$) from iteration policy learning (imitation of $\pi$).
% A compute-aware variant adds a per-iteration cost or a stricter $\tau$; a brief end-to-end fine-tune is optional but not required for strong performance.
